\section{Experimental Setup}\label{sec:method}

\subsection{Platforms, Sessions, and Software}\label{sec:platforms}
\begin{table*}[t]
\caption{\capTabPlatforms}\label{tab:platforms}
\centering\small\gentable{tab_platforms}\gennote{tab_platforms}
\end{table*}

A platform is one of our three hardware targets, which are the TPU v5e chip of the course cluster, the NVIDIA T4 GPU on Google Colab, and the Google Colab CPU runtime with two logical CPUs of an Intel Xeon processor. % src: recomputed: SEP.cpu.0915.logical_cpus, SEP.cpu.0915.cpu_model
A session is one Colab runtime or one Kubernetes Job in which a set of runs took place.
Table~\ref{tab:platforms} lists the stand-alone baseline run of each session.

The TPU runs used a Google Kubernetes Engine cluster of the course, in the us-west4 region. Each Kubernetes Job
requested the whole slice of eight chips. Some Jobs ran several experiments.
The Job definitions in the repository start from a \texttt{python:3.12-slim} image and install \texttt{jax[tpu]==0.10.2}. % src: configs/af_spike_job.yaml ("image: python:3.12-slim"); recomputed: V2.jax_pin
The Job definitions pin no other package and set a lower bound only for \texttt{tensorflow-cpu}.
The Job definitions clone AlphaFold2 from its default branch, and the runs did not record the commit of AlphaFold2.
 The TPU runs did not record their dates.
Most TPU results are in the first commit of our benchmark repository. This commit is dated August 8, 2026, 17:23 UTC. % src: recomputed: V2.first_commit
We committed the results of the matched grid (Section~\ref{sec:res-pmap}) on August 8, 2026, at 18:05 UTC. We committed the results of the external ensemble on August 11, 2026, at 22:38 UTC. % src: recomputed: V2.grid_commit, V2.ensemble_commit

The baseline runs on the T4 and on the CPU runtime took place in two Colab sessions on August 8, 2026, at 04:50 UTC and 05:19 UTC, respectively. % src: recomputed: AUG.gpu.first_epoch_utc, AUG.cpu.first_epoch_utc
These runs used a copy of the benchmark script inside the notebooks. We call this copy the notebook copy.
The notebooks in commit 2755250 hold this copy and the output of the runs. This copy differs from the scripts in \texttt{src/}. % src: recomputed: V2.first_commit; git show 2755250:notebooks/alphafold_cpu_benchmark.ipynb, cell 8 (script_content, output "Wrote 7449 bytes")
The first calls of these two Colab runs used the profiler. Table~\ref{tab:platforms} does not show these first calls.
The two Colab sessions of August 8, 2026, did not record their JAX version or their host CPU. % src: recomputed: AUG.gpu.first_epoch_utc, AUG.cpu.first_epoch_utc
The file \texttt{results/comparison.md} of the repository on August 12, 2026 (commit 612c5a2, 17:37 UTC) describes the CPU runtime as Colab's standard allocation of an Intel Xeon with two vCPUs. % src: recomputed: V4D.commit_612c5a2, V4D.comparison_612c5a2_xeon_2vcpu; the README gained it in commit e22f7af)
We cannot identify the script of the baseline run on the TPU.
The result file of the TPU baseline run lacks the field \texttt{host\_processor}, which the notebook copy writes. The result file also contains none of the extra fields that the committed script writes. % src: results/result_tpu-v5e-podslice.json -> keys; results/result_cpu-colab.json -> keys; recomputed: B.tpu.has_host_processor_field
The format of the file is closer to the notebook copy than to the committed script.

From September 15 to 17, 2026, we ran the committed script on the T4 and on the CPU runtime. % src: recomputed: V2.sep_dates
The CPU runtime ran in three sessions, on September 15, 2026, at 18:04 UTC, on September 16, 2026, at 12:37 UTC, and on September 17, 2026, at 08:52 UTC. % src: recomputed: V4C.sep_cpu_0915_start, V4C.sep_cpu_0916_start, V4C.sep_cpu_0917_start
The T4 ran in one session, on September 15, 2026, at 19:01 UTC. Table~\ref{tab:platforms} explains these times. % src: recomputed: V4C.sep_gpu_0915_start
In the three CPU sessions, the CPU runtime had an Intel Xeon processor at 2.20~GHz with two logical CPUs. % src: recomputed: SEP.cpu.0915.cpu_model, SEP.cpu.0916.cpu_model, SEP.cpu.0917.cpu_model, SEP.cpu.0915.logical_cpus, SEP.cpu.0916.logical_cpus, SEP.cpu.0917.logical_cpus
Each session ran the script twice. One run used the profiler, and the other did not. For these four sessions, tables
and text report the runs without the profiler. Section~\ref{sec:res-first} also uses the runs with the profiler. Each
of the four sessions recorded its commit, its package versions, and its host CPU.
These sessions ran commits 86fca03 (September 15, 2026, 17:57 UTC) and 9e1a07e (September 16, 2026, 12:35 UTC), which have the same benchmark scripts. % src: recomputed: V2.sep_src_identical, V4C.commit_86fca03, V4C.commit_9e1a07e; results/repro/2026-09-15_gpu-t4/environment_gpu-t4-repro.json -> repo_commit; results/repro/2026-09-16_cpu-colab/environment_cpu-colab-repro.json -> repo_commit
All four September sessions ran JAX 0.10.2 and AlphaFold2 at commit c77e5d2~\cite{alphafoldrepo}. % src: recomputed: SEP.cpu.0915.jax, SEP.cpu.0916.jax, SEP.cpu.0917.jax, SEP.gpu.0915.jax, SEP.cpu.0915.alphafold_commit
We did not repeat the TPU runs. Our access to the cluster ended with the course.
% The plan of the repeat is in commit 1651c42 (18 September 2026).
In the folder \texttt{cloud/}, we prepared a repeat of the TPU runs on Cloud TPU. We had no cloud resources when we
prepared the repeat, and we did not run it later.

\subsection{Workload and Inputs}\label{sec:workload}
\begin{table*}[t]
\caption{\capTabExperiments}\label{tab:experiments}
\centering\small\setlength{\tabcolsep}{3.5pt}\gentable{tab_experiments}\gennote{tab_experiments}
\end{table*}

Every run, except those of the model comparison, used \texttt{model\_3} of AlphaFold2.
This model needs no templates. Templates
are known structures of related proteins that \texttt{model\_1} and \texttt{model\_2} take as input. % src: AlphaFold2 at commit c77e5d2, alphafold/model/config.py:56-70 (template settings only for model_1 and model_2)
The script built the model with random weights. Haiku creates the random weights in the parameter initialization call.
 The
script never loaded the trained checkpoint. The MSA held only the query sequence, and no run searched a sequence
database. The featurization pads the MSA features to a fixed number of rows.
With one sequence, \texttt{model\_3} still runs on 512 MSA rows and 5120 extra MSA rows~\cite{alphafoldrepo}. % src: AlphaFold2 at commit c77e5d2, alphafold/model/config.py:233 ("'max_msa_clusters': 512"), config.py:73 ("'data.common.max_extra_msa': 5120"), alphafold/model/tf/data_transforms.py:389-425 (make_fixed_size); paper-v4/SOURCE_NOTES.md section 3
Only the first MSA row holds the query, and the other rows are padding.
Every run used zero recycles, one ensemble member, float32 parameters, and random seed 0, unless an experiment varied one of these settings. % src: src/spike_tpu_forward_pass.py ("random_seed=0")
Seed 0 sets both the parameter initialization and the featurization. In the external ensemble, only the featurization seed varies, from 0 to 7. % src: src/spike_ensemble_shard_forward_pass.py ("random_seed=i", "random_seed=0"); recomputed: V2.ensemble_seeds
AlphaFold2's source ran unmodified. Newer JAX releases renamed the arguments \texttt{a\_min} and \texttt{a\_max} of
\texttt{jax.numpy.clip} to \texttt{min} and \texttt{max}. A small wrapper in each script renamed these arguments.

The experiments used two input families. The fixed sequence has 118 residues. % src: recomputed: INPUT.toy_length
The repeated-alphabet sequence cycles through the 20 amino-acid letters up to the target length. % src: recomputed: V2.alphabet_length, INPUT.per_experiment ("alphabet ACDEFGHIKLMNPQRSTVWY repeated to the length")
In runs with several queries, each query starts one letter later in the alphabet. Table~\ref{tab:experiments} gives
the family of each experiment. The baseline configuration is the single-query script (Section~\ref{sec:experiments})
on the fixed sequence with \texttt{model\_3}, zero recycles, and float32 parameters.
The baseline configuration times \texttt{predict}.

\subsection{Timing, Memory, and Metrics}\label{sec:timing}
The scripts time each call of the model with \texttt{time.time} on the host. They wait for the outputs of the first
call and of the steady-state calls with \texttt{jax.block\_until\_ready}. Parameter initialization is the
\texttt{init\_params} call, which traces the model and creates the random weights. The first call is the earliest call
of \texttt{predict} or \texttt{apply} at a new input shape. The first call traces, compiles, and executes the model. A
steady-state call is a later call at the same input shape. A steady-state call reuses the compiled program.
 The single-query
and automatic-partitioning scripts time the parameter initialization, the first call, and the steady-state calls.
The
\texttt{vmap}, \texttt{pmap}, and ensemble scripts time the first call and one steady-state call. The Colab runs of August 8, 2026,
and all TPU runs timed one steady-state call, and the runs of September 15 to 17, 2026, without the profiler timed five. % src: src/spike_tpu_forward_pass.py ("--num_steady_state_runs=5"); recomputed: AUG.gpu.first_epoch_utc, V2.sep_dates

The scripts time two different regions. The single-query script and the automatic-partitioning script time
\texttt{predict}, which includes the confidence metrics on the host. The \texttt{vmap}, \texttt{pmap}, and ensemble
scripts time \texttt{apply} (Table~\ref{tab:experiments}).
The ensemble script runs \texttt{jax.lax.pmean} inside the same \texttt{jax.pmap} call as \texttt{apply}, so the timed region of the ensemble script also contains this collective operation. % src: src/spike_ensemble_shard_forward_pass.py ("logits_mean = jax.lax.pmean", "pmapped_apply = jax.pmap(sharded_apply")
 Except in the runs of September 15 to 17, 2026, without the profiler, the
single-query script and the \texttt{vmap} script run the first call inside \texttt{jax.profiler.trace}. This context
records a profile of the call. The profiler teardown is the time that the \texttt{jax.profiler.trace} context takes to
close after the first call returns. The teardown lies inside the timer of the first call. % src: recomputed: V2.sep_dates

Throughput is the number of queries that steady-state calls complete per second. We report throughput in queries/s.
The speedup is the throughput on $n$ chips divided by the throughput on one chip, with the same script, input family,
sequence length, and timed region. With multi-query \texttt{pmap}, each chip runs one query, so the speedup is a
weak-scaling measure of throughput and not a shorter time for one query. The one-chip reference of this speedup is
\texttt{pmap} on one chip. The parallel efficiency is the speedup divided by $n$. On one chip, multi-query
\texttt{pmap} runs one query. Every \texttt{pmap} run had all eight chips visible. We did not measure the utilization
of a chip. In this paper, a chip that works is a chip that holds data and runs a replica.

% Every script in src/ records stats.get("bytes_in_use") of device.memory_stats() for each
% device of jax.devices(). The definition of MB is in the note of Table III (tools/make_tables.py).
HBM in use is the value of \texttt{bytes\_in\_use} that JAX reports for a chip after the first steady-state call. This
value is the device memory allocated right after the first steady-state call. Peak memory is the value of
\texttt{peak\_bytes\_in\_use} that JAX reports for a chip at the same moment. We did not check whether this value
includes the temporary buffers of compiled programs.

\subsection{Experiments}\label{sec:experiments}
Table~\ref{tab:experiments} lists the baseline configuration and the 12 TPU experiments, with script, input, timed region, and number of runs. % src: recomputed: COUNT.tpu_experiments_with_raw_data
Experiments in one Job ran one after another in the same container, in a fixed order. Most settings ran once. The
repeated-runs experiment ran the baseline configuration three times on the TPU, in one Job.
Across these three runs, the coefficient of variation was 0.12\% for the steady-state call, 0.82\% for the parameter initialization, and 1.00\% for the first call. % src: results/sweep/repeated_runs.json -> runs; recomputed: RR.steady_state_seconds.cv_pct, RR.init_params_seconds.cv_pct, RR.first_predict_seconds.cv_pct
We computed the coefficient of variation with the sample standard deviation.
The steady-state times are rounded to 1~ms, about 0.2\% of the call, so the resolution of the records limits this coefficient. % src: src/spike_tpu_forward_pass.py ("round(t_steady_state, 3)"); recomputed: V4.ms_over_call_pct ("1 ms")
We ran the baseline configuration 11 times on the TPU. Ten of these runs started with an empty compilation cache. % src: recomputed: TPU.single_query_steady_n, V2.tpu_first_n
The note of Table~\ref{tab:experiments} lists the settings that make up these runs. We also tried to make only two or
four chips visible to the single-query script. The file \texttt{results/sweep/chip\_visibility.md} records that the
runs with two or four visible chips stopped with an error. We have no log of these runs.

\subsection{Provenance of the Result Files}\label{sec:records}
Most TPU results were committed together with their scripts. The files in \texttt{results/sweep/} are summaries of the
Job output. These summaries are not the files that the scripts wrote. Most field names in these summaries differ from
those that the scripts write.
Every memory value in the summaries is a whole number of megabytes. % src: recomputed: CV.all_bytes_multiple_of_1e6
The logs of the Jobs are not in the repository, and we have not found these logs. The first author produced these
summaries from the output of the Jobs with Claude (Anthropic). We have no record of how each value was transferred. We
report the values in these summaries unchanged. We checked the summaries against each other. In the \texttt{vmap} and
\texttt{pmap} summaries, each throughput equals the number of queries divided by the duration of the steady-state call,
within the rounding of that duration.
The 11 single-query runs of the baseline configuration come from seven files. These runs took 0.469~s to 0.471~s. % src: recomputed: TPU.single_query_steady_n, TPU.single_query_steady_min_max; paper-v4/SOURCE_NOTES.md ("seven files")
These checks cannot replace the output of the Jobs.
